% Part: second-order-logic
% Chapter: metatheory
% Section: second-order-arithmetic

\documentclass[../../../include/open-logic-section]{subfiles}

\begin{document}

\olfileid{sol}{met}{spa}

\olsection{द्वितीय-क्रम अंकगणित}

%READERNOTE{SOL6-A094: विगमन रूपबंधातील इतर मुक्त parameters सार्विकरीत्या बद्ध घेतले आहेत. Order proof मध्ये tail चा अंतर्भाव ही त्या सूत्राच्या पूर्ततेसाठी आवश्यक अट आहे; ती पुरेशी होण्यासाठी successor closure हवी. Tail स्वतः वरील सूत्र पूर्ण करतो, म्हणून ते पूर्ण करणाऱ्या सर्व संचांचा छेद योग्य order relation देतो.}
पेआनो अंकगणिताची उपपत्ती~$\Th{PA}$ हिच्यात
$\Th{Q}$ ची आठ स्वयंसिद्धके समाविष्ट आहेत, हे आठवा:
\begin{align*}
& \lforall[x][\eq/[x'][\Obj 0]]\\
& \lforall[x][\lforall[y][(\eq[x'][y'] \lif \eq[x][y])]]\\
& \lforall[x][(\eq[x][\Obj 0] \lor \lexists[y][\eq[x][y']])]\\
& \lforall[x][\eq[(x + \Obj 0)][x]]\\
& \lforall[x][\lforall[y][\eq[(x + y')][(x + y)']]]\\
& \lforall[x][\eq[(x \times \Obj 0)][\Obj 0]]\\
& \lforall[x][\lforall[y][\eq[(x \times y')][((x \times y) + x)]]]\\
& \lforall[x][\lforall[y][(x < y \liff \lexists[z][\eq[(z' + x)][y]])]]\\
\intertext{आणि पुढील स्वरूपाची सर्व वाक्येही:}
& (!A(\Obj 0) \land \lforall[x][(!A(x) \lif !A(x'))]) \lif \lforall[x][!A(x)].
\end{align*}
शेवटचे हे एक ``रूपबंध'' आहे: अंकगणिताच्या भाषेतील
अमर्याद !!{sentence}s निर्माण करणारा नमुना, प्रत्येक
!!{formula}~$!A(x)$ साठी एक वाक्य. इतर मुक्त चल असतील,
तर आरंभी त्यांच्यावर सार्विक संख्यापक जोडून वाक्यरूप घेतो. याला (प्रथम-क्रम)
\emph{विगमनाचा स्वयंसिद्धक-रूपबंध} म्हणतो.
\emph{द्वितीय-क्रम} पेआनो अंकगणित~$\Th{PA^2}$ मध्ये
विगमन एकाच वाक्यात मांडता येते. $\Th{PA^2}$ मध्ये
वरील पहिली आठ स्वयंसिद्धके आणि (द्वितीय-क्रम)
\emph{विगमनाचे स्वयंसिद्धक} असते:
\[
\lforall[X][((X(\Obj 0) \land \lforall[x][(X(x) \lif X(x'))]) \lif \lforall[x][X(x)])].
\]
याचा अर्थ असा: !!{domain} चा उपसंच~$X$ यामध्ये
$\Assign{\Obj{0}}{M}$ असेल आणि कोणत्याही~$x \in \Domain{M}$
बरोबर~$\Assign{\prime}{M}(x)$ ही असेल (म्हणजे तो
``उत्तराधिकारी क्रियेखाली बंद'' असेल), तर !!{domain} मधील
सर्व घटक त्यात असतील (म्हणजे $X =
\Domain{M})$.

विगमनाचे स्वयंसिद्धक असा निर्वाळा देते की ते पूर्ण करणाऱ्या
कोणत्याही !!{structure} च्या~$\Domain{M}$ मध्ये
स्वयंसिद्धकांना आवश्यक असलेलेच !!{element}s असतात;
म्हणजे $n \in \Nat$ साठी~$\num{n}$ ची मूल्ये.
$\Th{PA^2}$ च्या प्रतिमानात अमानक संख्या नसतात.

\begin{thm}
\ollabel{thm:sol-pa-standard}
जर $\Sat{M}{\Th{PA^2}}$ असेल, तर $\Domain{M} =
\Setabs{\Value{\num{n}}{M}}{n \in \Nat}$.
\end{thm}

\begin{proof}
$N = \Setabs{\Value{\num{n}}{M}}{n \in \Nat}$ असे मानू
आणि $\Sat{M}{\Th{PA^2}}$ असे गृहीत धरू. अर्थात,
कोणत्याही $n \in \Nat$ साठी
$\Value{\num{n}}{M} \in \Domain{M}$; म्हणून
$N \subseteq \Domain{M}$.

आता उलट दिशेचा अंतर्भाव दाखवू. $s(X) = N$ असे
चल-मूल्यांकन $s$ घेऊ. गृहीतकानुसार,
\begin{align*}
\Sat{M}{\lforall[X][((X(\Obj 0) \land \lforall[x][(X(x)
      \lif X(x'))]) \lif \lforall[x][X(x)])]}, & \text{ म्हणून}\\
\Sat{M}{(X(\Obj 0) \land \lforall[x][(X(x) \lif X(x'))]) \lif
  \lforall[x][X(x)]}[s]. &
\end{align*}
या सशर्त विधानाचे पूर्वांग पाहू. $\Value{\Obj{0}}{M} \in
N$, म्हणून $\Sat{M}{X(\Obj{0})}[s]$. दुसरे संयुक्तपद,
$\lforall[x][(X(x) \lif X(x'))]$, हेही पूर्ण होते.
कारण $x \in N$ असे समजू. $N$ च्या व्याख्येनुसार,
काही~$n$ साठी $x = \Value{\num{n}}{M}$. त्यावरून
$\Assign{\prime}{M}(x) = \Value{\num{n+1}}{M} \in N$.
म्हणून $\Assign{\prime}{M}(x) \in N$.

त्यामुळे $\Sat{M}{X(\Obj 0) \land \lforall[x][(X(x) \lif X(x'))]}[s]$.
म्हणून $\Sat{M}{\lforall[x][X(x)]}[s]$. याचा अर्थ,
प्रत्येक $x \in \Domain{M}$ साठी $x \in s(X) = N$.
म्हणून $\Domain{M} \subseteq N$.
\end{proof}

\begin{cor}
\ollabel{cor:sol-pa-categorical}
$\Th{PA^2}$ ची कोणतीही दोन प्रतिमाने समरूप असतात.
\end{cor}

\begin{proof}
\olref{thm:sol-pa-standard} नुसार,
$\Th{PA^2}$ च्या कोणत्याही प्रतिमानाच्या प्रांतात
$\Value{\num{n}}{M}$ ही मूल्येच असतात. असे प्रत्येक
प्रतिमान~$\Th{Q}$ चेही प्रतिमान असते.
\olref[mod][mar][stm]{prop:thq-standard} नुसार
ते मानक असते, म्हणजे~$\Struct{N}$ शी समरूप असते.
\end{proof}

वर $\Th{PA^2}$ ची व्याख्या पहिली आठ अंकगणितीय स्वयंसिद्धके
आणि द्वितीय-क्रम विगमनाचे स्वयंसिद्धक असलेली उपपत्ती म्हणून
केली. प्रत्यक्षात, द्वितीय-क्रम तर्कशास्त्राच्या
अभिव्यक्तीक्षमतेमुळे द्वितीय-क्रम पेआनो अंकगणितासाठी
अंकगणितीय स्वयंसिद्धकांपैकी फक्त \emph{पहिली दोन}
आणि विगमन पुरेसे असते.

\begin{prop}\ollabel{prop:sol-pa-definable}
पहिली दोन अंकगणितीय स्वयंसिद्धके (उत्तराधिकारी स्वयंसिद्धके)
आणि द्वितीय-क्रम विगमनाचे स्वयंसिद्धक असलेली द्वितीय-क्रम
उपपत्ती $\Th{PA^{2\dagger}}$ असो. मग $\le$, $+$ आणि
$\times$ हे $\Th{PA^{2\dagger}}$ मध्ये परिभाष्य आहेत.
\end{prop}

\begin{proof}
$\le$ परिभाष्य आहे हे दाखवण्यासाठी असे
सूत्र~$!A_\le(x, y)$ शोधावे लागेल की
$\Sat{N}{!A_\le(\num n, \num m)}$ जर आणि फक्त जर
$n \le m$. पुढील सूत्र पाहू:
\[
!B(x, Y) \ident Y(x) \land \lforall[y][(Y(y) \lif Y(y'))]
\]
संच $Y \subseteq \Nat$ कडून $!B(\num n, Y)$ ची
पूर्तता झाल्यास $\Setabs{m}{n \le m} \subseteq Y$ असते.
उलट, फक्त हा अंतर्भाव पुरेसा नाही; Y उत्तराधिकारी क्रियेखाली
बंदही असावा लागतो. n पासून पुढच्या सर्व नैसर्गिक संख्यांचा
संच स्वतः वरील सूत्राची पूर्तता करतो आणि अशा प्रत्येक Y मध्ये
तो समाविष्ट असतो. म्हणून
$!A_\le(x, y) \ident
\lforall[Y][(!B(x, Y) \lif Y(y))]$ असे घेता येते.

बेरीज परिभाष्य आहे हे पाहण्यासाठी लक्षात घ्या की $k+l=m$
जर आणि फक्त जर असे फलन $u$ असेल की $u(0)=k$,
प्रत्येक $n$ साठी $u(n')=u(n)'$ आणि $m = u(l)$.
या समतुल्यतेने $\Th{PA^{2\dagger}}$ मध्ये पुढील
सूत्राद्वारे बेरीज परिभाषित करता येते:
\[
!A_+(x,y,z) \ident \lexists[u][(u(\Obj 0)=x \land \lforall[w][u(w')=u
 (w)'] \land u(y)=z)]
\]
स्पष्टच, $\Sat{N}{!A_+(\num k, \num l, \num m)}$
जर आणि फक्त जर $k+l=m$.
\end{proof}

\begin{prob}
\olref[sol][met][spa]{prop:sol-pa-definable} चा पुरावा पूर्ण करा.
\end{prob}

\end{document}
